% Options for packages loaded elsewhere
\PassOptionsToPackage{unicode}{hyperref}
\PassOptionsToPackage{hyphens}{url}
\PassOptionsToPackage{dvipsnames,svgnames*,x11names*}{xcolor}
%
\documentclass[
  10pt,
]{article}
\usepackage{lmodern}
\usepackage{amssymb,amsmath}
\usepackage{ifxetex,ifluatex}
\ifnum 0\ifxetex 1\fi\ifluatex 1\fi=0 % if pdftex
  \usepackage[T1]{fontenc}
  \usepackage[utf8]{inputenc}
  \usepackage{textcomp} % provide euro and other symbols
\else % if luatex or xetex
  \usepackage{unicode-math}
  \defaultfontfeatures{Scale=MatchLowercase}
  \defaultfontfeatures[\rmfamily]{Ligatures=TeX,Scale=1}
  \setmainfont[]{DejaVu Serif}
  \setmonofont[]{DejaVu Sans Mono}
\fi
% Use upquote if available, for straight quotes in verbatim environments
\IfFileExists{upquote.sty}{\usepackage{upquote}}{}
\IfFileExists{microtype.sty}{% use microtype if available
  \usepackage[]{microtype}
  \UseMicrotypeSet[protrusion]{basicmath} % disable protrusion for tt fonts
}{}
\makeatletter
\@ifundefined{KOMAClassName}{% if non-KOMA class
  \IfFileExists{parskip.sty}{%
    \usepackage{parskip}
  }{% else
    \setlength{\parindent}{0pt}
    \setlength{\parskip}{6pt plus 2pt minus 1pt}}
}{% if KOMA class
  \KOMAoptions{parskip=half}}
\makeatother
\usepackage{xcolor}
\IfFileExists{xurl.sty}{\usepackage{xurl}}{} % add URL line breaks if available
\IfFileExists{bookmark.sty}{\usepackage{bookmark}}{\usepackage{hyperref}}
\hypersetup{
  colorlinks=true,
  linkcolor=Maroon,
  filecolor=Maroon,
  citecolor=Blue,
  urlcolor=Blue,
  pdfcreator={LaTeX via pandoc}}
\urlstyle{same} % disable monospaced font for URLs
\usepackage[margin=0.6in,landscape]{geometry}
\usepackage{longtable,booktabs}
% Correct order of tables after \paragraph or \subparagraph
\usepackage{etoolbox}
\makeatletter
\patchcmd\longtable{\par}{\if@noskipsec\mbox{}\fi\par}{}{}
\makeatother
% Allow footnotes in longtable head/foot
\IfFileExists{footnotehyper.sty}{\usepackage{footnotehyper}}{\usepackage{footnote}}
\makesavenoteenv{longtable}
\setlength{\emergencystretch}{3em} % prevent overfull lines
\providecommand{\tightlist}{%
  \setlength{\itemsep}{0pt}\setlength{\parskip}{0pt}}
\setcounter{secnumdepth}{-\maxdimen} % remove section numbering
\renewcommand{\arraystretch}{1.6}

\author{}
\date{}

\begin{document}

\hypertarget{phase-0.3-triangle-solving-drills}{%
\section{Phase 0.3 --- Triangle-Solving
Drills}\label{phase-0.3-triangle-solving-drills}}

\textbf{Type:} Tool drill, no mission narrative. \textbf{Target time:}
\textasciitilde30 min. \textbf{Prerequisite:} 0.0 Slide rule basics, 0.1
Trig/log table interpolation.

Much of hand orbital mechanics --- tracking geometry, Gauss's method,
range- and-bearing problems --- reduces at some point to careful
triangle solving. This drill covers the three tools you'll lean on
repeatedly: law of cosines, law of sines, and the ambiguous SSA case
(which trips up more people than it should, and is worth meeting here
rather than mid-problem later).

\hypertarget{law-of-cosines}{%
\subsection{Law of cosines}\label{law-of-cosines}}

\begin{verbatim}
c² = a² + b² − 2ab·cos(C)          (SAS: two sides + included angle → third side)
cos(C) = (a² + b² − c²) / (2ab)    (SSS: three sides → any angle)
\end{verbatim}

\hypertarget{law-of-sines}{%
\subsection{Law of sines}\label{law-of-sines}}

\begin{verbatim}
a/sin(A) = b/sin(B) = c/sin(C)
\end{verbatim}

\hypertarget{the-ambiguous-case-ssa}{%
\subsection{The ambiguous case (SSA)}\label{the-ambiguous-case-ssa}}

Given two sides and a non-included angle, there can be \textbf{zero,
one, or two} valid triangles. If the given angle is acute and the side
opposite it is shorter than the other given side, check both the arcsine
solution and its supplement (180° − that angle) --- both may produce a
valid triangle. This is a genuine geometric ambiguity, not a mistake to
avoid; real tracking data resolved it with a second observation or
physical reasoning about which solution made sense.

\hypertarget{drill-1-sas}{%
\subsection{Drill 1 --- SAS}\label{drill-1-sas}}

Sides a = 42.3, b = 57.8, included angle C = 63.5°. Find side c.

\hypertarget{drill-2-sss}{%
\subsection{Drill 2 --- SSS}\label{drill-2-sss}}

Sides a = 35.0, b = 48.0, c = 61.0. Find angle C (opposite the longest
side).

\hypertarget{drill-3-asa}{%
\subsection{Drill 3 --- ASA}\label{drill-3-asa}}

Angle A = 41.0°, angle B = 67.0°, side a = 28.5 (opposite A). Find angle
C, side b, and side c.

\hypertarget{drill-4-ssa-ambiguous-case}{%
\subsection{Drill 4 --- SSA (ambiguous
case)}\label{drill-4-ssa-ambiguous-case}}

Side a = 22.0, side b = 30.0, angle A = 25.0° (opposite side a). Find
all valid solutions for angle B, and complete each triangle (angle C and
side c).

\begin{center}\rule{0.5\linewidth}{0.5pt}\end{center}

\hypertarget{worksheet}{%
\subsection{Worksheet}\label{worksheet}}

\begin{longtable}[]{@{}llll@{}}
\toprule
Drill & Given & Working & Answer\tabularnewline
\midrule
\endhead
1 & a=42.3, b=57.8, C=63.5° & & c =\tabularnewline
2 & a=35.0, b=48.0, c=61.0 & & C =\tabularnewline
3 & A=41.0°, B=67.0°, a=28.5 & & C=\_\_\_ b=\_\_\_
c=\_\_\_\tabularnewline
4a & a=22.0, b=30.0, A=25.0° & & B=\_\_\_ C=\_\_\_
c=\_\_\_\tabularnewline
4b & (second solution) & & B=\_\_\_ C=\_\_\_ c=\_\_\_\tabularnewline
\bottomrule
\end{longtable}

\end{document}
